\documentclass[11pt]{article}
\usepackage[utf8]{inputenc}
\usepackage[T1]{fontenc}
\usepackage{geometry}
\usepackage{xcolor}
\usepackage{booktabs}
\usepackage{array}
\usepackage{enumitem}
\usepackage{amsmath}
\usepackage{amssymb}
\usepackage{hyperref}
\usepackage{listings}
\usepackage{fancyvrb}

\geometry{margin=1in}

% Custom colors
\definecolor{codeblue}{RGB}{0, 102, 204}
\definecolor{codegray}{RGB}{128, 128, 128}
\definecolor{codegreen}{RGB}{0, 128, 0}
\definecolor{codepurple}{RGB}{128, 0, 128}

% Code styling
\lstdefinestyle{pythonstyle}{
    language=Python,
    basicstyle=\ttfamily\footnotesize,
    keywordstyle=\color{codeblue},
    commentstyle=\color{codegreen},
    stringstyle=\color{codepurple},
    numberstyle=\tiny\color{codegray},
    breaklines=true,
    captionpos=b,
    frame=single,
    numbers=left,
    showstringspaces=false
}

\title{\textbf{Multi-Modal Stock TFT Implementation Roadmap} \\ 
       \large{Detailed Technical Specifications and Code Structure}}
\author{Implementation Guide - August 2025}
\date{}

\begin{document}

\maketitle

\section{Current Codebase Analysis}

Based on the existing codebase analysis, the project has a solid foundation with the following components already implemented:

\subsection{Existing Infrastructure}
\begin{itemize}
    \item ✅ \textbf{TFT Core:} Production-ready TFT implementation in \texttt{tft\_multimodal.py}
    \item ✅ \textbf{Data Pipeline:} Robust multimodal data loading via \texttt{dataModule/}
    \item ✅ \textbf{Training:} Unified training pipeline in \texttt{unified\_tft\_pipeline.py}
    \item ✅ \textbf{Trading Simulation:} Kelly criterion, DCA strategies in \texttt{trading\_simulator.py}
    \item ✅ \textbf{Visualization:} OHLC plotting with predictions in \texttt{ohlc\_plotter.py}
    \item ✅ \textbf{Caching:} Intelligent data caching system
    \item ⚠️ \textbf{Baselines:} Partial implementation in \texttt{fixed\_baseline\_pipeline.py}
\end{itemize}

\subsection{Current Baseline Models}
\begin{itemize}
    \item ✅ Linear Regression
    \item ✅ Random Forest Regressor
    \item ✅ XGBoost (optional, if installed)
    \item ❌ LSTM variants
    \item ❌ Transformer baselines
    \item ❌ Financial time series models
    \item ❌ Deep learning ensemble methods
\end{itemize}

\section{Implementation Tasks by Priority}

\subsection{Priority 1: Complete Baseline Suite (Days 1-7)}

\subsubsection{Task 1.1: Extend \texttt{fixed\_baseline\_pipeline.py}}

Add the following model implementations to the existing baseline runner:

\begin{lstlisting}[style=pythonstyle, caption=Enhanced Baseline Models]
# Add to fixed_baseline_pipeline.py after line 398

# LSTM-based models
from tensorflow.keras.models import Sequential
from tensorflow.keras.layers import LSTM, Dense, Dropout, Bidirectional

class LSTMWrapper:
    """Keras LSTM wrapper for sklearn-like interface"""
    def __init__(self, sequence_length=60, units=50, layers=2):
        self.sequence_length = sequence_length
        self.units = units
        self.layers = layers
        self.model = None
        
    def fit(self, X, y):
        # Reshape X for LSTM: (samples, timesteps, features)
        X_lstm = X.reshape(-1, self.sequence_length, X.shape[1]//self.sequence_length)
        
        self.model = Sequential()
        self.model.add(LSTM(self.units, return_sequences=True, 
                           input_shape=(self.sequence_length, X.shape[1]//self.sequence_length)))
        self.model.add(Dropout(0.2))
        
        for _ in range(self.layers - 1):
            self.model.add(LSTM(self.units, return_sequences=True))
            self.model.add(Dropout(0.2))
            
        self.model.add(LSTM(self.units))
        self.model.add(Dropout(0.2))
        self.model.add(Dense(y.shape[1] if len(y.shape) > 1 else 1))
        
        self.model.compile(optimizer='adam', loss='mse')
        self.model.fit(X_lstm, y, epochs=50, batch_size=32, verbose=0)
        
    def predict(self, X):
        X_lstm = X.reshape(-1, self.sequence_length, X.shape[1]//self.sequence_length)
        return self.model.predict(X_lstm, verbose=0)

# Add ARIMA models
from statsmodels.tsa.arima.model import ARIMA
from statsmodels.tsa.statespace.sarimax import SARIMAX

class ARIMAWrapper:
    """ARIMA wrapper for multivariate prediction"""
    def __init__(self, order=(1,1,1)):
        self.order = order
        self.models = []
        
    def fit(self, X, y):
        self.models = []
        for i in range(y.shape[1] if len(y.shape) > 1 else 1):
            target = y[:, i] if len(y.shape) > 1 else y
            model = ARIMA(target, order=self.order)
            fitted_model = model.fit()
            self.models.append(fitted_model)
            
    def predict(self, X):
        predictions = []
        for model in self.models:
            pred = model.forecast(steps=1)
            predictions.append(pred[0] if hasattr(pred, '__len__') else pred)
        return np.array(predictions).reshape(1, -1)

# Update the base_models dictionary
enhanced_models = {
    'LSTM_Vanilla': LSTMWrapper(units=50, layers=1),
    'LSTM_Deep': LSTMWrapper(units=100, layers=3),
    'LSTM_Bidirectional': None,  # Implement separately
    'ARIMA': ARIMAWrapper(order=(2,1,2)),
    'SARIMAX': None,  # Seasonal ARIMA
}
\end{lstlisting}

\subsubsection{Task 1.2: Create Advanced Ensemble Methods}

Create \texttt{ensemble\_baselines.py}:

\begin{lstlisting}[style=pythonstyle, caption=Ensemble Baseline Methods]
#!/usr/bin/env python3
"""
Advanced Ensemble Methods for Stock Prediction Baselines
========================================================
"""

import numpy as np
from sklearn.ensemble import VotingRegressor, StackingRegressor
from sklearn.model_selection import cross_val_score
from sklearn.metrics import mean_squared_error

class AdaptiveEnsemble:
    """Dynamically weighted ensemble based on recent performance"""
    
    def __init__(self, base_models, window_size=252):  # 1 trading year
        self.base_models = base_models
        self.window_size = window_size
        self.weights = None
        self.performance_history = []
        
    def fit(self, X, y):
        # Train all base models
        for name, model in self.base_models.items():
            model.fit(X, y)
            
        # Initialize equal weights
        self.weights = {name: 1.0/len(self.base_models) 
                       for name in self.base_models.keys()}
                       
    def predict(self, X):
        predictions = {}
        for name, model in self.base_models.items():
            predictions[name] = model.predict(X)
            
        # Weighted average prediction
        final_pred = np.zeros_like(list(predictions.values())[0])
        for name, pred in predictions.items():
            final_pred += self.weights[name] * pred
            
        return final_pred
        
    def update_weights(self, y_true, predictions):
        """Update weights based on recent performance"""
        errors = {}
        for name, pred in predictions.items():
            errors[name] = mean_squared_error(y_true, pred)
            
        # Inverse error weighting
        total_inv_error = sum(1.0/error for error in errors.values())
        for name in self.weights.keys():
            self.weights[name] = (1.0/errors[name]) / total_inv_error

class MetaLearningEnsemble:
    """Meta-learning ensemble using model predictions as features"""
    
    def __init__(self, base_models, meta_model):
        self.base_models = base_models
        self.meta_model = meta_model
        
    def fit(self, X, y):
        # Train base models with cross-validation to get meta-features
        meta_features = []
        
        for name, model in self.base_models.items():
            # Cross-validation predictions as meta-features
            cv_predictions = cross_val_score(model, X, y, cv=5, 
                                           scoring='neg_mean_squared_error')
            meta_features.append(cv_predictions)
            
            # Train on full dataset
            model.fit(X, y)
            
        meta_X = np.column_stack(meta_features)
        self.meta_model.fit(meta_X, y)
        
    def predict(self, X):
        # Get base model predictions
        base_predictions = []
        for name, model in self.base_models.items():
            pred = model.predict(X)
            base_predictions.append(pred.flatten())
            
        meta_X = np.column_stack(base_predictions)
        return self.meta_model.predict(meta_X)
\end{lstlisting}

\subsection{Priority 2: Multi-Asset Evaluation Framework (Days 8-21)}

\subsubsection{Task 2.1: Create Stock Grouping System}

Create \texttt{stock\_groups.py}:

\begin{lstlisting}[style=pythonstyle, caption=Stock Categorization System]
#!/usr/bin/env python3
"""
Stock Grouping and Sector Analysis
==================================
"""

STOCK_GROUPS = {
    'technology': {
        'symbols': ['AAPL', 'MSFT', 'GOOGL', 'NVDA', 'AMZN', 'TSLA', 'META', 'NFLX'],
        'description': 'Large-cap technology stocks',
        'characteristics': ['High volatility', 'Growth-oriented', 'Innovation-driven']
    },
    'financial': {
        'symbols': ['JPM', 'BAC', 'WFC', 'GS', 'MS', 'C', 'BRK-B'],
        'description': 'Major financial institutions',
        'characteristics': ['Interest rate sensitive', 'Economic cycle dependent']
    },
    'healthcare': {
        'symbols': ['JNJ', 'PFE', 'UNH', 'MRCK', 'ABBV', 'TMO', 'DHR'],
        'description': 'Healthcare and pharmaceutical companies',
        'characteristics': ['Defensive', 'Regulatory sensitive', 'Stable cash flows']
    },
    'consumer_discretionary': {
        'symbols': ['AMZN', 'TSLA', 'HD', 'MCD', 'NKE', 'SBUX', 'DIS'],
        'description': 'Consumer discretionary companies',
        'characteristics': ['Economic cycle sensitive', 'Consumer spending dependent']
    },
    'consumer_staples': {
        'symbols': ['KO', 'PEP', 'WMT', 'PG', 'JNJ', 'PFE'],
        'description': 'Consumer staples and essentials',
        'characteristics': ['Defensive', 'Stable demand', 'Dividend-focused']
    },
    'energy': {
        'symbols': ['XOM', 'CVX', 'COP', 'EOG', 'SLB', 'MPC'],
        'description': 'Energy sector companies',
        'characteristics': ['Commodity sensitive', 'Cyclical', 'High volatility']
    },
    'indices': {
        'symbols': ['SPY', 'QQQ', 'DIA', 'IWM', 'VTI'],
        'description': 'Major market indices ETFs',
        'characteristics': ['Diversified', 'Market representative', 'Lower volatility']
    }
}

MARKET_CONDITIONS = {
    'bull_market_2020_2021': {
        'start_date': '2020-03-23',
        'end_date': '2021-11-05',
        'description': 'COVID recovery bull market',
        'characteristics': ['Low interest rates', 'Fiscal stimulus', 'Tech outperformance']
    },
    'bear_market_2022': {
        'start_date': '2022-01-01',
        'end_date': '2022-10-12',
        'description': 'Inflation and rate hike bear market',
        'characteristics': ['Rising rates', 'Inflation concerns', 'Value outperformance']
    },
    'volatile_2023': {
        'start_date': '2023-01-01',
        'end_date': '2023-12-31',
        'description': 'Banking crisis and AI boom volatility',
        'characteristics': ['Banking stress', 'AI hype', 'Selective outperformance']
    }
}

def get_comparison_pairs():
    """Return meaningful stock comparison pairs for analysis"""
    return [
        ('AAPL', 'MSFT'),      # Tech giants
        ('JPM', 'BAC'),        # Major banks
        ('KO', 'PEP'),         # Beverage rivals
        ('SPY', 'QQQ'),        # Broad vs tech-heavy market
        ('XOM', 'CVX'),        # Energy majors
    ]

def get_sector_representatives():
    """Return one representative stock per major sector"""
    return {
        'Technology': 'AAPL',
        'Financial': 'JPM', 
        'Healthcare': 'JNJ',
        'Consumer Discretionary': 'AMZN',
        'Consumer Staples': 'KO',
        'Energy': 'XOM',
        'Market Index': 'SPY'
    }
\end{lstlisting}

\subsubsection{Task 2.2: Create Comprehensive Evaluation Runner}

Create \texttt{comprehensive\_evaluation.py}:

\begin{lstlisting}[style=pythonstyle, caption=Multi-Asset Evaluation Framework]
#!/usr/bin/env python3
"""
Comprehensive Multi-Asset Evaluation Framework
=============================================
"""

import pandas as pd
import numpy as np
from typing import Dict, List, Tuple
from concurrent.futures import ThreadPoolExecutor, as_completed
import logging
from stock_groups import STOCK_GROUPS, MARKET_CONDITIONS

class ComprehensiveEvaluator:
    """Runs systematic evaluation across multiple assets and conditions"""
    
    def __init__(self, models: Dict, config: Dict):
        self.models = models
        self.config = config
        self.results = {}
        
    def run_full_evaluation(self):
        """Run complete evaluation across all stocks and conditions"""
        
        # 1. Individual stock evaluation
        self.evaluate_by_stock_groups()
        
        # 2. Market condition analysis
        self.evaluate_by_market_conditions()
        
        # 3. Cross-sectional analysis
        self.evaluate_cross_sectional()
        
        # 4. Pair trading analysis
        self.evaluate_pair_trading()
        
        return self.results
        
    def evaluate_by_stock_groups(self):
        """Evaluate models performance by stock sector/group"""
        
        for group_name, group_info in STOCK_GROUPS.items():
            print(f"\\n📊 Evaluating {group_name.upper()} sector...")
            
            group_results = {}
            symbols = group_info['symbols']
            
            # Parallel evaluation of stocks in group
            with ThreadPoolExecutor(max_workers=4) as executor:
                future_to_symbol = {
                    executor.submit(self.evaluate_single_stock, symbol): symbol 
                    for symbol in symbols
                }
                
                for future in as_completed(future_to_symbol):
                    symbol = future_to_symbol[future]
                    try:
                        result = future.result()
                        group_results[symbol] = result
                    except Exception as exc:
                        print(f'❌ {symbol} generated exception: {exc}')
                        
            self.results[f'group_{group_name}'] = group_results
            
    def evaluate_by_market_conditions(self):
        """Evaluate model performance under different market conditions"""
        
        for condition_name, condition_info in MARKET_CONDITIONS.items():
            print(f"\\n📈 Evaluating {condition_name} period...")
            
            # Filter data by date range
            start_date = condition_info['start_date']
            end_date = condition_info['end_date']
            
            condition_results = {}
            
            # Test representative stocks in this period
            test_symbols = ['AAPL', 'JPM', 'SPY']  # Representative cross-section
            
            for symbol in test_symbols:
                result = self.evaluate_single_stock(
                    symbol, 
                    start_date=start_date, 
                    end_date=end_date
                )
                condition_results[symbol] = result
                
            self.results[f'condition_{condition_name}'] = condition_results
            
    def evaluate_single_stock(self, symbol: str, start_date=None, end_date=None):
        """Evaluate all models on a single stock"""
        
        # Load data for symbol with date filters
        data_loader = self.get_data_loader(symbol, start_date, end_date)
        X_train, y_train, X_test, y_test = data_loader.get_train_test_split()
        
        stock_results = {}
        
        for model_name, model in self.models.items():
            try:
                # Train model
                model.fit(X_train, y_train)
                
                # Generate predictions
                y_pred = model.predict(X_test)
                
                # Calculate metrics
                metrics = self.calculate_metrics(y_test, y_pred)
                
                # Calculate trading performance
                trading_metrics = self.calculate_trading_metrics(
                    y_test, y_pred, symbol
                )
                
                stock_results[model_name] = {
                    'prediction_metrics': metrics,
                    'trading_metrics': trading_metrics
                }
                
            except Exception as e:
                print(f"❌ Error with {model_name} on {symbol}: {e}")
                stock_results[model_name] = {'error': str(e)}
                
        return stock_results
\end{lstlisting}

\subsection{Priority 3: Statistical Analysis Framework (Days 22-28)}

\subsubsection{Task 3.1: Create Statistical Testing Suite}

Create \texttt{statistical\_analysis.py}:

\begin{lstlisting}[style=pythonstyle, caption=Statistical Analysis Framework]
#!/usr/bin/env python3
"""
Statistical Analysis and Significance Testing
=============================================
"""

import numpy as np
import pandas as pd
from scipy import stats
from scipy.stats import friedmanchisquare
import scikit_posthocs as sp
from typing import Dict, List, Tuple

class StatisticalAnalyzer:
    """Comprehensive statistical analysis of model performance"""
    
    def __init__(self, results: Dict):
        self.results = results
        self.analysis_results = {}
        
    def run_full_analysis(self):
        """Run complete statistical analysis"""
        
        # 1. Descriptive statistics
        self.calculate_descriptive_stats()
        
        # 2. Significance testing
        self.run_significance_tests()
        
        # 3. Effect size analysis
        self.calculate_effect_sizes()
        
        # 4. Risk analysis
        self.analyze_risk_metrics()
        
        return self.analysis_results
        
    def calculate_descriptive_stats(self):
        """Calculate descriptive statistics for all models"""
        
        model_performance = self.aggregate_model_performance()
        
        desc_stats = {}
        for metric in ['mse', 'mae', 'sharpe_ratio', 'max_drawdown']:
            desc_stats[metric] = {}
            
            for model_name in model_performance.keys():
                values = model_performance[model_name][metric]
                
                desc_stats[metric][model_name] = {
                    'mean': np.mean(values),
                    'std': np.std(values),
                    'median': np.median(values),
                    'q25': np.percentile(values, 25),
                    'q75': np.percentile(values, 75),
                    'min': np.min(values),
                    'max': np.max(values)
                }
                
        self.analysis_results['descriptive_stats'] = desc_stats
        
    def run_significance_tests(self):
        """Run statistical significance tests comparing models"""
        
        model_performance = self.aggregate_model_performance()
        significance_results = {}
        
        for metric in ['mse', 'mae', 'sharpe_ratio']:
            print(f"\\n🧮 Testing significance for {metric}...")
            
            # Prepare data for Friedman test
            model_names = list(model_performance.keys())
            data_matrix = []
            
            for model_name in model_names:
                data_matrix.append(model_performance[model_name][metric])
                
            # Friedman test (non-parametric ANOVA for repeated measures)
            statistic, p_value = friedmanchisquare(*data_matrix)
            
            significance_results[metric] = {
                'friedman_statistic': statistic,
                'p_value': p_value,
                'significant': p_value < 0.05
            }
            
            # Post-hoc analysis if significant
            if p_value < 0.05:
                df = pd.DataFrame(data_matrix, index=model_names).T
                posthoc = sp.posthoc_nemenyi_friedman(df)
                significance_results[metric]['posthoc'] = posthoc
                
        self.analysis_results['significance_tests'] = significance_results
        
    def calculate_effect_sizes(self):
        """Calculate effect sizes for model comparisons"""
        
        model_performance = self.aggregate_model_performance()
        effect_sizes = {}
        
        # Cohen's d for pairwise comparisons
        model_names = list(model_performance.keys())
        
        for metric in ['mse', 'mae', 'sharpe_ratio']:
            effect_sizes[metric] = {}
            
            for i, model1 in enumerate(model_names):
                for j, model2 in enumerate(model_names[i+1:], i+1):
                    
                    values1 = model_performance[model1][metric]
                    values2 = model_performance[model2][metric]
                    
                    # Cohen's d
                    pooled_std = np.sqrt(((len(values1)-1)*np.var(values1) + 
                                        (len(values2)-1)*np.var(values2)) / 
                                       (len(values1) + len(values2) - 2))
                    
                    cohens_d = (np.mean(values1) - np.mean(values2)) / pooled_std
                    
                    effect_sizes[metric][f'{model1}_vs_{model2}'] = {
                        'cohens_d': cohens_d,
                        'magnitude': self.interpret_cohens_d(cohens_d)
                    }
                    
        self.analysis_results['effect_sizes'] = effect_sizes
        
    @staticmethod
    def interpret_cohens_d(d):
        """Interpret Cohen's d effect size"""
        abs_d = abs(d)
        if abs_d < 0.2:
            return 'negligible'
        elif abs_d < 0.5:
            return 'small'
        elif abs_d < 0.8:
            return 'medium'
        else:
            return 'large'
\end{lstlisting}

\subsection{Priority 4: Paper Generation Framework (Days 26-31)}

\subsubsection{Task 4.1: Automated Results Table Generation}

Create \texttt{paper\_generator.py}:

\begin{lstlisting}[style=pythonstyle, caption=Automated Paper Generation]
#!/usr/bin/env python3
"""
Automated Research Paper Content Generation
===========================================
"""

import pandas as pd
import matplotlib.pyplot as plt
import seaborn as sns
from typing import Dict, Any

class PaperGenerator:
    """Generate LaTeX tables and figures for research paper"""
    
    def __init__(self, results: Dict, analysis: Dict):
        self.results = results
        self.analysis = analysis
        
    def generate_main_results_table(self):
        """Generate main results comparison table"""
        
        latex_table = r"""
\\begin{table}[htbp]
\\centering
\\caption{Model Performance Comparison Across All Stocks}
\\label{tab:main_results}
\\begin{tabular}{lcccccc}
\\toprule
Model & MAE & RMSE & Dir. Acc. & Sharpe & Max DD & Win Rate \\\\
\\midrule
"""
        
        # Extract aggregated results
        for model_name in self.get_model_names():
            metrics = self.get_aggregated_metrics(model_name)
            
            latex_table += f"{model_name} & "
            latex_table += f"{metrics['mae']:.4f} & "
            latex_table += f"{metrics['rmse']:.4f} & "
            latex_table += f"{metrics['directional_accuracy']:.3f} & "
            latex_table += f"{metrics['sharpe_ratio']:.3f} & "
            latex_table += f"{metrics['max_drawdown']:.3f} & "
            latex_table += f"{metrics['win_rate']:.3f} \\\\\\\\"
            
        latex_table += r"""
\\bottomrule
\\end{tabular}
\\end{table}
"""
        
        return latex_table
        
    def generate_sector_analysis_table(self):
        """Generate sector-wise performance analysis"""
        
        latex_table = r"""
\\begin{table}[htbp]
\\centering
\\caption{TFT Performance by Market Sector}
\\label{tab:sector_analysis}
\\begin{tabular}{lccccc}
\\toprule
Sector & Sharpe Ratio & Max Drawdown & Total Return & Volatility & \\# Stocks \\\\
\\midrule
"""
        
        for sector in ['Technology', 'Financial', 'Healthcare', 'Consumer', 'Energy']:
            sector_metrics = self.get_sector_metrics(sector)
            
            latex_table += f"{sector} & "
            latex_table += f"{sector_metrics['sharpe_ratio']:.3f} & "
            latex_table += f"{sector_metrics['max_drawdown']:.3f} & "
            latex_table += f"{sector_metrics['total_return']:.3f} & "
            latex_table += f"{sector_metrics['volatility']:.3f} & "
            latex_table += f"{sector_metrics['num_stocks']} \\\\\\\\"
            
        latex_table += r"""
\\bottomrule
\\end{tabular}
\\end{table}
"""
        
        return latex_table
        
    def generate_significance_table(self):
        """Generate statistical significance results table"""
        
        significance = self.analysis['significance_tests']
        
        latex_table = r"""
\\begin{table}[htbp]
\\centering
\\caption{Statistical Significance Tests (Friedman Test)}
\\label{tab:significance}
\\begin{tabular}{lccc}
\\toprule
Metric & Test Statistic & p-value & Significant \\\\
\\midrule
"""
        
        for metric, results in significance.items():
            latex_table += f"{metric.upper()} & "
            latex_table += f"{results['friedman_statistic']:.3f} & "
            latex_table += f"{results['p_value']:.4f} & "
            latex_table += f"{'Yes' if results['significant'] else 'No'} \\\\\\\\"
            
        latex_table += r"""
\\bottomrule
\\end{tabular}
\\end{table}
"""
        
        return latex_table
        
    def generate_all_tables(self):
        """Generate all tables for the paper"""
        
        tables = {
            'main_results': self.generate_main_results_table(),
            'sector_analysis': self.generate_sector_analysis_table(),
            'significance': self.generate_significance_table()
        }
        
        return tables
\end{lstlisting}

\section{Daily Task Breakdown}

\subsection{Week 1: Baseline Implementation}

\textbf{Day 1-2:} LSTM and Deep Learning Baselines
\begin{itemize}
    \item Implement LSTM variants (vanilla, bidirectional, stacked)
    \item Add GRU models and CNN-LSTM hybrid
    \item Integrate with existing \texttt{fixed\_baseline\_pipeline.py}
\end{itemize}

\textbf{Day 3-4:} Traditional ML Enhancement
\begin{itemize}
    \item Add LightGBM, CatBoost, Extra Trees
    \item Implement SVM with different kernels
    \item Add ensemble methods (voting, stacking)
\end{itemize}

\textbf{Day 5-6:} Financial Time Series Models
\begin{itemize}
    \item Implement ARIMA, GARCH, VAR models
    \item Add Markov switching models
    \item Create naive baselines (random walk, moving averages)
\end{itemize}

\textbf{Day 7:} Testing and Integration
\begin{itemize}
    \item Test all baseline models
    \item Fix integration issues
    \item Validate output consistency
\end{itemize}

\subsection{Week 2-3: Multi-Asset Evaluation}

\textbf{Day 8-10:} Stock Group Analysis
\begin{itemize}
    \item Implement \texttt{stock\_groups.py}
    \item Run technology sector analysis
    \item Run financial sector analysis
\end{itemize}

\textbf{Day 11-14:} Comprehensive Stock Coverage
\begin{itemize}
    \item Healthcare and consumer sectors
    \item Energy and utilities sectors
    \item Index ETF analysis
\end{itemize}

\textbf{Day 15-17:} Market Condition Analysis
\begin{itemize}
    \item Bull market periods (2020-2021)
    \item Bear market periods (2022)
    \item Volatile periods (2023-2024)
\end{itemize}

\textbf{Day 18-21:} Cross-sectional Analysis
\begin{itemize}
    \item Pair trading analysis
    \item Correlation studies
    \item Portfolio-level evaluation
\end{itemize}

\subsection{Week 4: Statistical Analysis}

\textbf{Day 22-24:} Core Statistical Testing
\begin{itemize}
    \item Implement \texttt{statistical\_analysis.py}
    \item Run Friedman tests for model comparison
    \item Calculate effect sizes (Cohen's d)
\end{itemize}

\textbf{Day 25-26:} Advanced Analytics
\begin{itemize}
    \item Risk-adjusted performance metrics
    \item Regime-dependent analysis
    \item Robustness testing
\end{itemize}

\textbf{Day 27-28:} Results Compilation
\begin{itemize}
    \item Aggregate all results
    \item Generate summary statistics
    \item Create visualization pipeline
\end{itemize}

\subsection{Week 5: Paper Writing}

\textbf{Day 29-30:} Content Generation
\begin{itemize}
    \item Implement \texttt{paper\_generator.py}
    \item Generate all LaTeX tables
    \item Create figures and plots
\end{itemize}

\textbf{Day 31:} Final Assembly
\begin{itemize}
    \item Compile complete paper
    \item Final review and editing
    \item Submission preparation
\end{itemize}

\section{Success Metrics}

\subsection{Technical Deliverables}
\begin{itemize}
    \item ✅ 15+ baseline models implemented and tested
    \item ✅ 25+ stocks analyzed across 6 sectors
    \item ✅ 3+ market conditions evaluated
    \item ✅ Statistical significance established
    \item ✅ Complete paper with tables and figures
\end{itemize}

\subsection{Research Quality}
\begin{itemize}
    \item ✅ Rigorous statistical methodology
    \item ✅ Comprehensive baseline comparison
    \item ✅ Multi-dimensional evaluation
    \item ✅ Reproducible results
    \item ✅ Clear contribution statement
\end{itemize}

\vfill
\begin{center}
\textit{Implementation roadmap created: July 17, 2025}
\end{center}

\end{document}
